\item \subquestionpoints{3}
All this while you should have stayed away from the test data completely. Now that
you have convinced yourself that the model is working as expected (i.e., the
observations you made in the previous part matches what you learnt in class
about regularization), it is finally time to measure the model performance on
the test set. Once we measure the test set performance, we report it (whatever
value it may be), and NOT go back and refine the model any further.

Initialize your model from the parameters saved in part (a) (i.e., the non-regularized
model), and evaluate the model performance on the test data. Repeat this using the
parameters saved in part (b) (i.e., the regularized model).

Report your test accuracy for both regularized model and non-regularized model.  
Briefly (in one sentence) explain why this outcome makes sense.
You should get test accuracy close to $0.9287$ without regularization and $0.9676$ with
regularization. Small deviations from these numbers are fine (they can depend on details such as
the random initialization); in all cases you should observe good accuracy overall and \emph{higher}
test accuracy with regularization.